\chapter{Conclusion}
\label{ch:conclusion}

This project built and validated a numerical framework for computing charged particle orbits in prescribed electromagnetic fields, then applied it to study particle dynamics in the magnetic dipole field.

The solver was verified against analytic solutions for gyromotion, helical motion, $\mathbf{E} \times \mathbf{B}$ drift, and gradient-$B$ drift, with energy conservation below $10^{-6}$ in every test. Applied to the dipole field, it demonstrated bounce motion, confirmed
approximate conservation of $\mu$ and showed that the guiding-centre integrator reproduces the slow bounce dynamics faithfully in the adiabatic regime. A simulation in physical SI units at $L = 3$ matchedthe analytic mirror latitude to within 0.02\,\%.

The framework was extended to tilted dipole fields at $47^\circ$ and $59^\circ$, reproducing field geometries representing Neptune and Uranus and showing that the apparent asymmetry in geographic coordinates is a projection effect. Adding a corotation electric field recovered the input rotation rate to within 96.6\,\%, and a rotating tilted dipole simulation combined all three adiabatic motions with time-dependent fields in a single integration.

The modular design means that adding a new field configuration requires only a new function in \texttt{fields.py}, with no changes to the solver.

\section*{Future Work}
Several extensions would build on this framework.
\begin{itemize}
  \item \textbf{Relativistic velocities.} Extending the equations of
    motion to include the relativistic Lorentz factor would allow
    modelling of higher-energy radiation belt particles, where
    relativistic corrections become significant above approximately
    100\,keV for electrons.
  \item \textbf{Higher-order drifts.} The guiding-centre expansion can
    be carried to higher order to include effects such as polarisation
    drift, which becomes relevant for time-varying electric fields and
    for particles with larger gyroradii.
  \item \textbf{Higher adiabatic invariants.} This project focused on
    the first adiabatic invariant $\mu$. Tracking the second invariant
    $J = \oint v_\parallel \, \mathrm{d}s$ and the third invariant
    $\Phi$ (the drift-shell flux) would allow a more complete test of
    adiabatic theory and connect directly to radiation belt dynamics.
  \item \textbf{Symplectic integrator.} Replacing RK45 with a symplectic
    scheme such as the Boris algorithm would suppress gradual phase
    drift, enabling accurate tracking over thousands of bounce periods.
\end{itemize}